%%%%%%%%%%%%%%%%%%%%%%%%%%%%%%%%%%%%%%%%%%%%%%%%%%%%%%%%%%%%%%%%%%%%%%%%%%%%%%%
% PROPUESTA DEL PROYECTO DE GRADO
%%%%%%%%%%%%%%%%%%%%%%%%%%%%%%%%%%%%%%%%%%%%%%%%%%%%%%%%%%%%%%%%%%%%%%%%%%%%%%%

\chapter{Propuesta del proyecto de grado}
\label{chap:propuesta}

La propuesta presenta la estructura prevista para el documento final y la
planificación de las actividades necesarias para cumplir los objetivos
específicos. Las duraciones y fechas del cronograma deberán establecerse con
base en la disponibilidad real del estudiante, del taller y de los componentes.

\section{Índice tentativo}
\label{sec:indice_tentativo}

El índice tentativo proyecta la organización del documento final. La estructura
se ajustará a los resultados obtenidos y a las observaciones del docente y del
tutor, manteniendo la organización institucional del perfil.

\begin{enumerate}
    \item \textbf{CAPÍTULO I: INTRODUCCIÓN}
    \begin{enumerate}
        \item Introducción
        \item Planteamiento del problema
        \begin{enumerate}
            \item Antecedentes del problema
            \item Formulación del problema
            \item Árbol de problemas
        \end{enumerate}
        \item Objetivos
        \begin{enumerate}
            \item Objetivo general
            \item Objetivos específicos
        \end{enumerate}
        \item Motivación
        \item Justificación
        \item Delimitación
        \begin{enumerate}
            \item Límites
            \item Alcances
        \end{enumerate}
    \end{enumerate}

    \item \textbf{CAPÍTULO II: MARCO TEÓRICO}
    \begin{enumerate}
        \item Estado del arte
        \begin{enumerate}
            \item Soluciones académicas para dosificación de materiales pulverulentos
            \item Soluciones comerciales de dosificación de harina
            \item Análisis comparativo de las soluciones revisadas
        \end{enumerate}
        \item Fundamentos teóricos
        \begin{enumerate}
            \item Dosificación gravimétrica
            \item Propiedades relevantes de la harina
            \item Sistemas mecánicos de alimentación de sólidos
            \item Instrumentación para pesaje
            \item Sistemas de accionamiento y control
            \item Criterios de evaluación del desempeño
        \end{enumerate}
        \item Conclusiones del capítulo
    \end{enumerate}

    \item \textbf{CAPÍTULO III: MARCO PRÁCTICO Y DESARROLLO DEL PROTOTIPO}
    \begin{enumerate}
        \item Metodología de desarrollo
        \item Requerimientos funcionales y técnicos
        \item Diseño del prototipo
        \begin{enumerate}
            \item Diseño mecánico
            \item Sistema de medición gravimétrica
            \item Sistema de control
            \item Integración de subsistemas
        \end{enumerate}
        \item Construcción e implementación
        \item Pruebas de funcionamiento y verificación de requerimientos
        \item Evaluación comparativa con el procedimiento manual
        \item Herramientas y materiales
        \item Conclusiones del capítulo
    \end{enumerate}

    \item \textbf{CAPÍTULO IV: RESULTADOS Y ANÁLISIS}
    \begin{enumerate}
        \item Resultados de las pruebas de funcionamiento
        \item Comparación del prototipo y el procedimiento manual
        \item Análisis de costos de construcción e integración
        \item Discusión de resultados
        \item Conclusiones del capítulo
    \end{enumerate}

    \item \textbf{CAPÍTULO V: CONCLUSIONES Y RECOMENDACIONES}
    \begin{enumerate}
        \item Conclusiones
        \item Recomendaciones
        \item Trabajos futuros
    \end{enumerate}
    \item Bibliografía
    \item Anexos
\end{enumerate}

\section{Cronograma de actividades}
\label{sec:cronograma_actividades}

El cronograma se organizará según los objetivos específicos. Cada actividad
deberá registrar su duración estimada en días, fecha de inicio y fecha de
finalización. Las tareas de redacción se incluyen explícitamente y se vinculan
con los capítulos correspondientes. Las fechas y duraciones quedan pendientes
de definición para evitar asignar plazos sin conocer la disponibilidad efectiva
ni el calendario acordado para el proyecto.

\begin{table}[htbp]
    \centering
    \caption{Actividades y plazos requeridos para elaborar el cronograma.}
    \label{tab:actividades_cronograma}
    \small
    \setlength{\tabcolsep}{4pt}
    \renewcommand{\arraystretch}{1.25}
    \begin{tabularx}{\textwidth}{@{}p{0.07\textwidth}X p{0.30\textwidth} c c c@{}}
        \toprule
        \textbf{OE} & \textbf{Actividad} & \textbf{Entregable / redacción} &
        \textbf{Días} & \textbf{Inicio} & \textbf{Fin} \\
        \midrule
        1 & Completar y analizar el diagnóstico de la dosificación manual; establecer los requerimientos funcionales y técnicos. & Requerimientos documentados & Por definir & Por definir & Por definir \\
        1 & Redactar y revisar los capítulos I y II conforme al diagnóstico, los requerimientos y la bibliografía consultada. & Redacción de Introducción y Marco teórico & Por definir & Por definir & Por definir \\
        2 & Desarrollar el diseño mecánico, el sistema de medición gravimétrica y el sistema de control conforme a los requerimientos. & Diseños y documentación técnica & Por definir & Por definir & Por definir \\
        2 & Redactar y revisar el diseño del prototipo, los cálculos y la integración prevista. & Redacción del diseño en el capítulo III & Por definir & Por definir & Por definir \\
        3 & Fabricar, ensamblar e integrar los subsistemas del prototipo. & Prototipo funcional integrado & Por definir & Por definir & Por definir \\
        3 & Redactar y revisar el proceso de construcción e implementación. & Redacción de implementación en el capítulo III & Por definir & Por definir & Por definir \\
        4 & Ejecutar pruebas de funcionamiento y registrar resultados para verificar los requerimientos. & Registros y verificación documentada & Por definir & Por definir & Por definir \\
        4 & Redactar los procedimientos de prueba y los resultados de verificación. & Redacción de resultados de pruebas & Por definir & Por definir & Por definir \\
        5 & Realizar ensayos comparables del prototipo y del procedimiento manual y analizar las variables de desempeño. & Resultados de comparación & Por definir & Por definir & Por definir \\
        5 & Redactar la comparación, discusión y análisis de resultados. & Redacción del capítulo IV & Por definir & Por definir & Por definir \\
        6 & Documentar materiales, componentes, fabricación y montaje; consolidar el costo de construcción e integración. & Desglose de costos & Por definir & Por definir & Por definir \\
        6 & Redactar conclusiones y recomendaciones; revisar coherencia, referencias, formato y versión final del documento. & Redacción del capítulo V y revisión final & Por definir & Por definir & Por definir \\
        \bottomrule
    \end{tabularx}
\end{table}

\subsection{Diagrama de Gantt}
\label{subsec:gantt}

El diagrama de Gantt se completará cuando se definan las fechas y duraciones
indicadas en la Tabla~\ref{tab:actividades_cronograma}. Cada barra deberá
representar la actividad correspondiente y mostrar la secuencia de trabajo por
objetivo específico, incluidas las tareas de redacción.

% Pendiente de incorporar el diagrama con fechas acordadas y duraciones en días.
